\notepage{53}{The hydrogen atom: radial quantization}
Use the Coulomb convention of the notes,
$V(r)=-e^2/r$. In SI units, each Coulomb factor $e^2$ here means
$e_{\rm SI}^2/(4\pi\epsilon_0)$. With an infinitely heavy nucleus,
\[
 U''+\frac{2m}{\hbar^2}
 \left[E+\frac{e^2}{r}-\frac{\ell(\ell+1)\hbar^2}{2mr^2}\right]U=0.
\]
For a bound state $E<0$, set
$\lambda=\sqrt{-2mE}/\hbar$. The limiting equations give
\begin{align*}
 r\to\infty:&\quad U''-\lambda^2U\simeq0,
       &&U\propto\ee^{-\lambda r},\\
 r\to0:&\quad U''-\frac{\ell(\ell+1)}{r^2}U\simeq0,
       &&U\propto r^{\ell+1}
       \quad\text{(regular solution)}.
\end{align*}
For $\ell=0$, the regular boundary condition $U(0)=0$ chooses the
same behavior. Introduce
\[
 \rho=\lambda r,\qquad U(\rho)=\ee^{-\rho}v(\rho),\qquad
 q=\frac{2me^2}{\hbar^2\lambda}.
\]
The differential equation is
\[
 v''-2v'+\left[\frac q\rho-\frac{\ell(\ell+1)}{\rho^2}\right]v=0.
\]
Expand $v=\rho^{\ell+1}\sum_{k=0}^{\infty}c_k\rho^k$. Substitution gives
\[
 c_{k+1}=\frac{2(k+\ell+1)-q}{(k+1)(k+2\ell+2)}c_k.
\]
Normalizable bound-state solutions require termination at some
$K=0,1,2,\ldots$, so $q=2(K+\ell+1)$. Therefore
\[
 E=-\frac{me^4}{2\hbar^2(K+\ell+1)^2}.
\]
\editorial{The rescaling coefficient has been written explicitly as
$q=2me^2/(\hbar^2\lambda)$ to keep the dimensionless equation
consistent with $\lambda$ as an inverse length.}
Define the principal quantum number $n=K+\ell+1$. Then
\[
 \boxed{E_n=-\frac{me^4}{2\hbar^2n^2}},\qquad
 n=1,2,\ldots,\quad \ell=0,1,\ldots,n-1.
\]
Each $\ell$ has $2\ell+1$ magnetic states. The orbital degeneracy is
$\sum_{\ell=0}^{n-1}(2\ell+1)=n^2$, before including spin.
